% =====================================================================
%  1 Projektübersicht  (Aufgabenstellung 1.1: „Beschreibung des Projekts
%  und des angestrebten Ergebnisses“)
%  Leitfragen Prüfungsleitfaden Phase 1:
%   - Was ist das Ziel meiner Arbeit?  - Was ist meine Idee?
%   - Ideenprüfung: Was spricht für, was gegen die Idee?
%   - Was ist das Konzept?  - Wie plane ich vorzugehen?
%   - Welche Methodik/welches Tool habe ich gewählt und warum?
%  Bewertung „Problemabgrenzung/Zielsetzung“ (10 %): Ziele, Umfang,
%  angestrebtes Ergebnis klar; Nutzen für die Zielgruppe deutlich.
%  Hinweis: Kein „ich“, „wir“, „man“ (Richtlinien Kap. 1).
%  Jedes (Unter-)Kapitel mind. ½ Seite; Unterteilung nur mit ≥ 2 Unterpunkten.
% =====================================================================
\section{Projektübersicht}
\TODO{2--3 Sätze Hinführung zu den Unterkapiteln}

\subsection{Ausgangssituation und Motivation}
\TODO{Problem, das die Work-Life-Balance-App löst; Kontext; Relevanz mit Quellen belegen}

\subsection{Zielsetzung und Kundennutzen}
\TODO{Ziel des Projekts; Zielgruppe (Details im Anforderungsdokument); klar definierter Customer Value}

\subsection{Projektumfang und Abgrenzung}
\TODO{Was ist Teil des Projekts (Kernfunktionen), was ausdrücklich nicht? Abgrenzung zu bestehenden Apps (Kreativität, Neuheit); Ideenprüfung: Was spricht dafür, was dagegen?}

\subsection{Angestrebtes Ergebnis}
\TODO{Funktionsfähige Windows-Desktopanwendung in Python mit GUI und Testdaten, Dokumentation, GitHub-Repository}

\subsection{Vorgehensmodell und Werkzeuge}
\TODO{Gewähltes Vorgehen (z. B. iterativ-inkrementell) mit Begründung und Quelle; Werkzeuge: LaTeX, draw.io, Git/GitHub, Python}
